\section{Related work} 

\textbf{Parameter settings related to weight decay.} 
\citet{Loshchilov2017DecoupledWD} explicitly highlights the coupling between the learning rate $\eta$ and the weight decay coefficient $\lambda$ in their study of Adam. For another parameter, \citet{blake2025umupunitscaledmaximalupdate} observe that, under fixed batch size and weight decay, the optimal learning rate decreases as the dataset size $N$ increases. This finding is broadly consistent with the results of \citet{wang2025setadamwsweightdecay}, who shows that larger datasets require a longer effective time scale $T$, which can be achieved by either reducing the weight decay or decreasing the learning rate. \citet{bergsma2025powerlinesscalinglaws} further demonstrates that the optimal weight decay coefficient scales linearly with the batch size $B$. Compared to directly tuning the learning rate, adjusting $\lambda$ provides improved stability in practice. In addition, \citet{fan2025robustlayerwisescalingrules} and \citet{wang2025setadamwsweightdecay} have shown that weight decay is critical to transfer efficacy. Similar to the conclusions given by \citet{wang2025setadamwsweightdecay}, \citet{dangelo2024needweightdecaymodern} observe that under under-training conditions, the optimal generalization performance of SGD often approximates $\eta \propto \frac{1}{\lambda}$. 

The work most similar to ours is \cite{wang2025setadamwsweightdecay}. They investigate the parameter settings of the AdamW algorithm by one-to-one mapping from the EMA timescale to the weight decay hyperparameter. However, in this work, we aim to study the coupling effects between $\lambda$ and other optimization parameters by deriving explicit expressions of $\lambda$ under generalization-driven analysis.

Due to space constraints, we have included the details of the related work in the \textbf{Appendix \ref{related}}.
